\subsection{Simple Regression Using Cement}

Simple linear regression was performed using both Python's
\texttt{statsmodels} package and ScalaTion. Both implementations produced
essentially identical coefficient estimates. The fitted model was

\[
\widehat{\mathrm{Strength}}
=
13.4425 + 0.07958(\mathrm{Cement}),
\]

where cement content is measured in kg/m$^3$ and compressive strength is
measured in MPa.

The positive slope indicates that higher cement content is associated with
greater compressive strength. A 1 kg/m$^3$ increase in cement content is
associated with an estimated increase of approximately 0.0796 MPa in
compressive strength. Equivalently, a 10 kg/m$^3$ increase in cement content
corresponds to an estimated increase of approximately 0.796 MPa.

The slope was statistically significant ($p<0.001$). ScalaTion reported
$R^2=0.247837$ and adjusted $R^2=0.247105$, consistent with the statsmodels
results. Thus, cement alone explained approximately 24.8\% of the observed
variation in concrete compressive strength. ScalaTion reported an RMSE of
approximately 14.48 MPa.

\subsection{Simple Regression Using Superplasticizer}

Simple linear regression was also performed using superplasticizer content as
the predictor. ScalaTion and statsmodels again produced essentially identical
results. The fitted equation was

\[
\widehat{\mathrm{Strength}}
=
29.4660 + 1.02373(\mathrm{Superplasticizer}),
\]

where superplasticizer content is measured in kg/m$^3$.

The positive slope indicates that greater superplasticizer content is
associated with greater compressive strength. A 1 kg/m$^3$ increase in
superplasticizer content is associated with an estimated increase of
approximately 1.024 MPa in compressive strength. The slope was statistically
significant ($p<0.001$).

ScalaTion reported $R^2=0.134014$ and adjusted $R^2=0.133171$, matching the
statsmodels results to rounding precision. Therefore, superplasticizer alone
explained approximately 13.4\% of the observed variation in compressive
strength. ScalaTion reported an RMSE of approximately 15.54 MPa.

\subsection{Regression Discussion}

Both selected predictors exhibited positive and statistically significant
linear relationships with concrete compressive strength, and ScalaTion and
statsmodels produced essentially identical results. Cement provided the
stronger simple regression model, with $R^2 \approx 0.248$, compared with
$R^2 \approx 0.134$ for superplasticizer. This finding is consistent with the
correlation analysis, in which cement also had the stronger Pearson
correlation with compressive strength.

Although both predictors were statistically significant, neither simple
regression explained the majority of the variation in compressive strength.
This suggests that concrete strength depends on multiple mixture components
and curing conditions rather than on either cement or superplasticizer alone.